\chapter{Brief Review of Partial Differential Equations}

%%%%%%%%%%%%%%%%%%%%%%%%%%%%%%%%
\section{First Order Equations}
%%%%%%%%%%%%%%%%%%%%%%%%%%%%%%%%

\subsection{Semi-explicit Solutions via the Method of Characteristics}

\paragraph{Advection Equation}
The solutions of the equation
\begin{equation}
    \label{eq:adveccion}
    \frac{\partial u}{\partial t} (t,x) + \bm v(t,x) \cdot \nabla u  = 0
\end{equation}
can be sought in the form of characteristics, that is, looking for curves $X_t$ such that $u(t,X_t) = u_0(X_0)$. Thus,
\begin{equation*}
    \frac{\partial u}{\partial t} (t,X_t) + \nabla u(t,X_t) \cdot \frac{dX_t}{dt} = 0.
\end{equation*}
It suffices to require that $X_t$ satisfy
\begin{equation}
    \label{eq:caracteristicas}
    \frac{dX_t}{dt} = \bm v(t, X_t).
\end{equation}
The Picard-Lindelöf theorem tells us that if $\bm v$ is globally Lipschitz in $x$ and continuous in $t$, then there exists a unique solution of \eqref{eq:caracteristicas} for a given $X_0$.

\medskip

For fixed $x$ we have curves $t \mapsto X_t(x)$ where $X_0(x) = x$. For fixed $t$, think of the transformation $X_t : x \mapsto X_t(x)$.
If the field $X_t$ is invertible for each time $t$ in a certain domain, the solution is of the form $u_t(t,x) = u_0 (X_t^{-1} (x))$.

\paragraph{Advection Equation with Zero Order Term}
In the equation
\begin{equation}
    \label{eq:adveccion con termino de orden cero}
    \frac{\partial u}{\partial t} (t,x) + \bm v(t,x) \cdot \nabla u + c(t,x) u = 0
\end{equation}
one can look for solutions of the form $A_t u(t,X_t) = u_0 (X_0)$. Using the chain rule again, we see that it is sufficient that $X_t$ satisfy \eqref{eq:caracteristicas} and that
\begin{equation*}
    \frac{dA_t}{dt} = A_t c(t, X_t), \qquad A_0 = 1.
\end{equation*}

\paragraph{Transport Equation}
The equation
\begin{equation}
    \frac{\partial u}{\partial t} (t,x) + \diver( \bm v(t,x) u) = 0
\end{equation}
can be written as \eqref{eq:adveccion con termino de orden cero} where $c = \diver \bm v$. This case is particularly important.

\paragraph{Burgers' Equation}

Nonlinear problems are more difficult. In $d=1$ we have Burgers' equation
\begin{equation}
    \label{eq:Burgers}
    \frac{\partial u}{\partial t} + u \frac{\partial u}{\partial x} = 0.
\end{equation}
In this equation $v = u$. But we observe that $v(t,X_t) = u(t,X_t) = u_0(X_0)$. Then the solution of \eqref{eq:caracteristicas} is precisely
\begin{equation*}
    X_t (x) = x + t u_0(x).
\end{equation*}
Another way to see this is that for classical solutions $u(t,y) = u_0(x)$ if the relation $y = x + t u_0(x)$ holds.
Let us see two extreme cases:
\begin{enumerate}
    \item If $u_0$ is non-decreasing, then $X_t : x \mapsto X_t (x)$ is strictly increasing in $x$, and therefore invertible for all $t > 0$. Given $y$ there exists a unique $x$ such that $y = x + t u_0(x)$. If $u_0$ is regular, then there exists a unique classical solution $u$ of \eqref{eq:Burgers}.

    \item Let $u_0(x) = -x$.
          For time $t \in (0,1)$ we have $X_t(x) = (t-1)x$ so $u(t,x) = u_0( \frac{x}{t-1} ) = \frac{x}{1-t}$. Clearly this solution blows up at $t = 1$.
          Let's see this through the characteristics.
          For time $t=1$ we have $X_1(x) = 0$ for any $x$. All characteristics passing through points of the form $(t,x) = (0,x)$ intersect at the point $(t,x) = (1,0)$. At this point, for any $x$
          \[
              u(1,0) = u_0(x) = -x.
          \]
          This is clearly a contradiction. Therefore, there is no classical solution defined at time $t = 1$.
\end{enumerate}

The method of characteristics admits a very general presentation for nonlinear problems. See [EVANS, ...] \todo{Add citation}

\subsection{Boundary Conditions}

\section{Solution semi-group}


\section{Elliptic Equations}

\begin{equation}
    \label{eq:Laplace}
    -\Delta u + a u = f.
\end{equation}

\subsection{Problems in Bounded Domains}

\subsubsection{Fourier Series on an Interval}

\subsubsection{Fourier Series on a Domain}

\subsubsection{Weak Solutions in a Domain. Sobolev Spaces}

\subsubsection{Maximum Principle. Accretivity of the Laplacian in $L^\infty$}

\subsubsection{Accretivity of the Laplacian in $L^1$}



\subsection{Problem in the Whole Space. Fourier Transform}


%%%%%%%%%%%%%%%%%%%%%%%%%%%%%%%%
\section{Parabolic Equations}
%%%%%%%%%%%%%%%%%%%%%%%%%%%%%%%%

\paragraph{The Heat Equation}
The canonical example of a linear parabolic equation is the heat equation
\begin{equation}
    \label{eq:calor}
    \frac{\partial u}{\partial t} -\Delta u = 0.
\end{equation}

\paragraph{The Viscous Burgers' Equation}
Usually, adding a Laplacian-type term makes the equations more regular. For example, one way to understand \eqref{eq:Burgers} is the Burgers' equation with a viscous term
\begin{equation}
    \label{eq:Burgers viscosa}
    \frac{\partial u}{\partial t} + u\frac{\partial u}{\partial x} = \varepsilon \frac{\partial^2 u}{\partial x^2}
\end{equation}

\subsubsection{Fourier Series in Domains}

\subsubsection{Semigroup Theory}

\subsubsection{$L^p$ Contractivity}

%%%%%%%%%%%%%%%%%%%%%%%%%%%%%%%%
\section{Hyperbolic Equations}
%%%%%%%%%%%%%%%%%%%%%%%%%%%%%%%%

\begin{equation}
    \tag{W}
    \label{eq:wave}
    \frac{\partial^2 u}{\partial t^2} - c^2\Delta u = 0.
\end{equation}

\subsubsection{Case $d=1$. d'Alembert's Formula}

If $u(x,0) = u_0(x)$ and $\frac{\partial u}{\partial t} (x,0) = v_0(x)$ then
\begin{equation*}
    \tag{W-d'A}
    \label{eq:wave 1d D'Alambert}
    u(x,t) = \frac{u_0(x-ct) + u_0(x+ct)}{2} + \frac 1 {2c} \int_{x-ct}^{x+ct} v_0(s) ds
\end{equation*}

\subsubsection{Difficulties in $d>1$}
